%%%%%%%%%%%%%%%%%%%%%%%%%%%%%%%%%%%%%%%%%%%%%%%%%%%%%%%%%%%%%%%%%%%%%%
% Lớp 10 | Đại số tổ hợp | Hoán vị. Chỉnh hợp. Tổ hợp
% Dạng 0D8?2-5: Bài toán chọn đối tượng khác
% Số câu: 36  (sinh tự động bởi tach_cau.py; từ nay thầy sửa trực tiếp file này)
%%%%%%%%%%%%%%%%%%%%%%%%%%%%%%%%%%%%%%%%%%%%%%%%%%%%%%%%%%%%%%%%%%%%%%

%%%%%%%%%%%%%%%% TRẮC NGHIỆM 4 PHƯƠNG ÁN %%%%%%%%%%%%%%%%

% nguồn: MaHoaNH_TracNghiem_Lop10
\begin{ex}%[0D8N2-5]%[ma:2a80efdc]
	Có $10$  cái bút khác nhau và $8$ quyển sách giáo khoa khác nhau. Một bạn học sinh cần chọn $1$ cái bút và $1$ quyển sách. Hỏi bạn học sinh đó có bao nhiêu cách chọn?
	 \choice
		 {\True $80$}
		{$70$} 
		{$60$} 
		{$90$} 
	\loigiai{ 
		Số cách chọn $1$ cái bút có $10$ cách,  số cách chọn $1$ quyển sách có $8$ cách.\\
		 Vậy theo quy tắc nhân, số cách chọn $1$ cái bút và $1$ quyển sách là $10\cdot 8=80$ cách.
	} 
\end{ex}

% nguồn: MaHoaNH_TracNghiem_Lop10
\begin{ex}%[0D8H2-5]%[ma:0b2ff2d0]
	Có $3$ bức thư và $3$ con tem. Hỏi có bao nhiêu cách dán tem lên $3$ bức thư, biết mỗi bức thư chỉ  dán 1 con tem.
	\choice 
		 {$3$}
		{\True  $6$}
		{$1$} 
		 {$2!$}
	\loigiai{ 
		Dán $3$ con  tem lên $3$ bức thư có $3!=6$ cách.
	 }
\end{ex}

% nguồn: MaHoaNH_TracNghiem_Lop10
\begin{ex}%[0D8H2-5]%[ma:9d063c2e]
	Một hộp bánh trung thu có $4$ bánh dẻo gồm các loại nhân khác  nhau và $6$ bánh nướng có các loại nhân khác nhau. Hỏi có bao nhiêu cách lấy ra $6$ cái bánh để phát cho các em thiếu nhi?
	\choice 
		 {$240$}
		{$151\,200$} 
		 {$14\,200$}
		{\True $210$} 
	 \loigiai{
		Mỗi cách lấy ra $6$ cái bánh trong $10$ bánh là một tổ hợp chập $6$ của tập $10$ phần tử.\\ 
		Vậy có $\mathrm{C}_{10}^6=210$ cách lấy ra  $6$ bánh để phát cho các em thiếu nhi.
	} 
\end{ex}

% nguồn: MaHoaNH_TracNghiem_Lop10
\begin{ex}%[0D8H2-5]%[ma:60f8a3fa]
	Cho tập  $S$ gồm $n$ điểm phân biệt trên mặt phẳng sao cho không có $3$ điểm nào thẳng hàng. Tìm $n$ sao cho số tam giác có $3$ đỉnh lấy từ $3$ điểm thuộc $S$ gấp đôi số đoạn thẳng được nối từ $3$ điểm thuộc $S$.
	 \choice
		 {$n=7$}
		{$n=12$} 
		 {$n=9$}
		 {\True $n=11$}
	 \loigiai{
		Ta có $\mathrm{C}_n^3=3\mathrm{C}_n^2$ $(1)$  (với $n\ge 3$, $n\in \mathbb{N}$)
		\[\Leftrightarrow \dfrac{n!}{3!\left(n-3\right)!}=\dfrac{3.n!}{2!\left(n-2\right)!} \Leftrightarrow  \dfrac{1}{6}=\dfrac{3}{2(n-2)} \Leftrightarrow n=11.\]
	 }
\end{ex}

% nguồn: MaHoaNH_TracNghiem_Lop10
\begin{ex}%[0D8H2-5]%[ma:d33bd322]
	Một hộp chứa $6$ quả cầu trắng và $4$ quả cầu đen. Có bao nhiêu cách lấy $2$  quả cầu cùng màu từ hộp đó?
	 \choice
		{\True  $21$}
		 {$42$}
		{$10$} 
		{$24$} 
	\loigiai{ 
		Trường hợp 1:  lấy 2 quả cầu trắng có $\mathrm{C}_6^2$ cách.\\
		Trường  hợp 2: lấy 2 quả cầu đen có $\mathrm{C}_4^2$ cách.\\
		Vậy có $\mathrm{C}_6^2+\mathrm{C}_4^2=21$  cách lấy thỏa mãn.
	 }
\end{ex}

% nguồn: MaHoaNH_TracNghiem_Lop10
\begin{ex}%[0D8H2-5]%[ma:0845c2ec]
	Có $7$ bông hoa khác nhau. Hỏi có bao nhiêu cách chọn $4$ bông hoa từ $7$ bông hoa đã cho để cắm vào $4$  lọ hoa khác nhau?
	 \choice
		{$35$} 
		{\True $840$} 
		 {$28$}
		{$24$} 
	\loigiai{ 
		Số cách chọn $4$ bông hoa từ $7$ bông hoa là  $\mathrm{C}_{7}^{4}$ (cách).\\
		Số cách cắm $4$ bông hoa đã chọn  vào $4$ lọ hoa là $4!$ (cách).\\
		Vậy có $\mathrm{C}_{7}^{4}\cdot 4!=840$ (cách). 
	} 
\end{ex}

% nguồn: MaHoaNH_TracNghiem_Lop10
\begin{ex}%[0D8H2-5]%[ma:74e004b1]
	Một hộp có $7$ quả cầu xanh, $9$ quả cầu đỏ và $12$ quả cầu vàng. Có bao nhiêu  cách chọn ra $3$ quả cầu có đúng $2$ màu?
	 \choice
		{$3\,276$} 
		 {$1\,095$}
		{$2\,859$} 
		{\True $2\,181$} 
	 \loigiai{
		 Số cách chọn $3$ quả cầu từ hộp là $\mathrm{C}_{28}^3=3\,276$ cách.\\
		 Số cách chọn $3$ quả cầu có đúng $1$ màu là $\mathrm{C}_7^3+\mathrm{C}_9^3+\mathrm{C}_{12}^3=339$ cách.\\
		Số cách chọn $3$ quả cầu  có đủ $3$ màu là $7\cdot 9\cdot 12=756$ cách.\\
		Khi đó, số cách chọn ra $3$ quả cầu có đúng  $2$ màu là $3\,276-339-756=2\,181$ cách.
	} 
\end{ex}

% nguồn: MaHoaNH_TracNghiem_Lop10
\begin{ex}%[0D8H2-5]%[ma:de0975e5]
	Một trường THPT có $10$ lớp $12$, mỗi lớp cử $3$ học sinh tham gia vẽ tranh cổ động. Các lớp tiến hành bắt tay giao lưu với nhau. Tính số lần bắt tay  của các học sinh với nhau, biết rằng hai học sinh khác nhau ở hai lớp khác nhau chỉ bắt tay đúng $1$ lần.
	 \choice
		{\True $405$} 
		{$435$} 
		 {$30$}
		{$45$} 
	\loigiai{ 
		Mỗi lớp cử ra $3$ học sinh nên $10$ lớp cử ra $30$  học sinh.\\
		Suy ra  số lần bắt tay là $\mathrm{C}_{30}^2$.\\
		Số lần  bắt tay của các học sinh học cùng một lớp là $10\cdot \mathrm{C}_3^2$.\\
		Vậy số lần bắt tay của các học sinh với nhau là  $\mathrm{C}_{30}^2-10\cdot \mathrm{C}_3^2=405$.
	 }
\end{ex}

% nguồn: MaHoaNH_TracNghiem_Lop10
\begin{ex}%[0D8H2-5]%[ma:284d9297]
	Một hộp chứa $18$ viên bi gồm $3$ bi đỏ, $6$ bi xanh và $9$ bi vàng. Chọn ngẫu nhiên từ hộp ra $2$ quả bi. Có  bao nhiêu cách để $2$ bi được chọn khác màu?
	 \choice
		{$72$} 
		{$45$} 
		{$153$} 
		 {\True $99$}
	 \loigiai{
		Chọn được $1$ bi đỏ  và $1$ bi xanh, có $\mathrm{C}_3^1\cdot \mathrm{C}_6^1=18$ cách.\\
		Chọn được $1$ bi đỏ và $1$ bi  vàng, có $\mathrm{C}_3^1\cdot \mathrm{C}_9^1=27$ cách.\\
		Chọn được $1$ bi vàng và $1$ bi xanh, có $\mathrm{C}_6^1\cdot \mathrm{C}_9^1=54$ cách.\\ 
		Vậy để chọn được $2$ bi khác màu có $18+27+54=99$  cách.
	 }
\end{ex}

% nguồn: MaHoaNH_TracNghiem_Lop10
\begin{ex}%[0D8H2-5]%[ma:90c1964e]
	Trên bàn có $8$ cây bút  chì khác nhau, $6$ cây bút bi khác nhau và $10$ cuốn tập khác nhau. Số cách khác nhau để chọn được đồng thời một cây bút chì, một cây bút bi và một cuốn tập.
	\choice 
		 {$48$}
		{$60$} 
		{$24$} 
		{\True  $480$}
	 \loigiai{
		\textbf{Công đoạn $1$:} Chọn $1$ cây bút chì từ $8$ cây bút chì có  $\mathrm{C}_{8}^1=8$ cách.\\
		\textbf{Công  đoạn $2$:} Chọn $1$ cây bút bi từ $6$ cây bút chì có $\mathrm{C}_{6}^1=6$ cách.\\
		\textbf{Công đoạn $3$:} Chọn $1$ cây  cuốn tập từ $10$ cuốn tập có $\mathrm{C}_{10}^1=10$ cách.\\
		Theo quy tắc  nhân có $8\cdot6\cdot10=480$ cách.
	} 
\end{ex}

% nguồn: MaHoaNH_TracNghiem_Lop10
\begin{ex}%[0D8H2-5]%[ma:4be77488]
	Trong kho đèn trang trí đang còn $6$ bóng đèn loại I, $8$ bóng đèn loại II, các bóng đèn đều khác nhau về màu sắc và hình dáng. Lấy ra $5$ bóng đèn bất kỳ. Hỏi có bao nhiêu khả năng xảy ra số bóng  đèn loại I nhiều hơn số bóng đèn loại II?
	 \choice
		{\True $686$} 
		 {$6\,480$}
		 {$680$}
		 {$6\,460$}
	\loigiai{ 
		Có $3$ trường hợp xảy ra.\\ 
		Trường  hợp 1: Lấy được $5$ bóng đèn loại I có $\mathrm{C}_6^5$ cách.\\
		Trường hợp 2: Lấy được $4$ bóng đèn loại I, $1$ bóng đèn loại II có $\mathrm{C}_6^4\cdot \mathrm{C}_8^1$  cách.\\
		Trường hợp 3: Lấy được $3$ bóng đèn loại I, $2$ bóng đèn loại II có  $\mathrm{C}_6^3\cdot \mathrm{C}_8^2$ cách.\\
		 Theo quy tắc cộng, có $\mathrm{C}_6^5+\mathrm{C}_6^4\cdot \mathrm{C}_8^1+\mathrm{C}_6^3\cdot \mathrm{C}_8^2=686$ cách .\\
	} 
\end{ex}

% nguồn: MaHoaNH_TracNghiem_Lop10
\begin{ex}%[0D8H2-5]%[ma:943f4d88]
	Với $4$ quyển sách toán khác nhau và $3$ quyển tập khác nhau, có bao nhiêu cách tặng một phần quà gồm một quyển sách toán và một  quyển tập cho $1$ em học sinh?
	\choice 
		{$4^3$ cách} 
		{$16$ cách} 
		{$7$  cách}
		 {\True $12$ cách}
	 \loigiai{
		\begin{itemize}
			\item Chọn một quyển sách  toán có $4$ cách.
			\item Chọn một quyển  tập có $3$ cách.
		\end{itemize}
		Theo quy tắc nhân ta có $4\cdot3=12$ cách tặng cho  $1$ em học sinh.
	 }
\end{ex}

% nguồn: MaHoaNH_TracNghiem_Lop10
\begin{ex}%[0D8H2-5]%[ma:36c92d09]
	Trong kho đèn trang trí đang còn  $5$ bóng đèn loại $I$, $7$ bóng đèn loại $II$, các bóng đèn đều khác nhau về màu sắc và hình dáng. Lấy ra $5$ bóng đèn bất kỳ. Hỏi có bao nhiêu khả năng xảy ra số bóng đèn loại $I$ nhiều hơn số bóng đèn loại $II$?
	\choice 
		{\True  $246$}
		 {$3360$}
		 {$3480$}
		 {$245$}
	\loigiai{ 
		Có  $3$ trường hợp xảy ra:
		\begin{enumerate}[TH 1:]
			\item Lấy được $5$ bóng  đèn loại $I$ có $1$ cách.
			\item Lấy  được $4$ bóng đèn loại $I$, $1$ bóng đèn loại $II$ có $\mathrm{C}_5^4\cdot \mathrm{C}_7^1$ cách.
			\item Lấy được $3$ bóng  đèn loại $I$, $2$ bóng đèn loại $II$ có $\mathrm{C}_5^3\cdot \mathrm{C}_7^2$ cách.
		\end{enumerate}
		Theo quy tắc cộng, có  $1+\mathrm{C}_5^4\cdot \mathrm{C}_7^1+\mathrm{C}_5^3\cdot \mathrm{C}_7^2=246$ cách.
	 }
\end{ex}

% nguồn: MaHoaNH_TracNghiem_Lop10
\begin{ex}%[0D8V2-5]%[ma:2be03cd3]
	Một  hộp đựng $20$ quả cầu được đánh số từ $1$ đến $20$. Lấy ba quả cầu từ hộp trên rồi cộng số ghi trên đó lại. Hỏi có tất cả bao nhiêu cách lấy để kết quả thu được là một số chia hết cho $3$?
	\choice 
		 {$390$}
		 {$120$}
		{\True $384$} 
		 {$502$}
	\loigiai{ 
		Chia $20$ quả cầu trong  hộp làm ba nhóm.\\
		Nhóm $1$ gồm các quả cầu mang số chia hết cho $3$, có $6$ quả.\\ 
		Nhóm $2$ gồm các quả cầu mang số chia cho $3$  dư $1$, có $7$ quả.\\
		Nhóm $3$ gồm các quả cầu mang số chia cho $3$ dư $2$, có  $7$ quả.\\
		Lấy ba quả cầu từ hộp  trên rồi cộng số ghi trên đó lại, được một số chia hết cho $3$ có các trường hợp sau\\
		Trường  hợp $1$: lấy $3$ quả cầu ở nhóm $1$, có $\mathrm{C}_6^3$ cách.\\
		Trường  hợp $2$: lấy $3$ quả cầu ở nhóm $2$, có $\mathrm{C}_7^3$ cách.\\
		Trường hợp $3$: lấy $3$ quả cầu ở nhóm $3$, có $\mathrm{C}_7^3$ cách.\\ 
		Trường hợp $4$: lấy mỗi phần $1$ quả cầu có $\mathrm{C}_6^1\cdot \mathrm{C}_7^1\cdot  \mathrm{C}_7^1$ cách.\\
		Vậy có $\mathrm{C}_6^3+\mathrm{C}_7^3+\mathrm{C}_7^3+\mathrm{C}_6^1\cdot \mathrm{C}_7^1\cdot  \mathrm{C}_7^1=384$ cách lấy được ba quả cầu thỏa mãn yêu cầu bài toán.
	 }
\end{ex}

% nguồn: MaHoaNH_TracNghiem_Lop10
\begin{ex}%[0D8V2-5]%[ma:7778122b]
	Thầy Tuấn có $10$ phần thưởng  khác nhau trong đó có $4$ cặp sách, $3$ hộp bút, $3$ quyển vở. Thầy muốn lấy ra $5$ món quà và tặng cho $5$ em học sinh $A$, $B$, $C$, $D$, $E$ mỗi em một phần thưởng. Hỏi thầy giáo có bao nhiêu cách tặng cho các em học sinh sao cho sau khi tặng xong, mỗi một trong ba loại phần thưởng trên đều còn ít nhất một phần thưởng.
	 \choice
		{$204$ cách} 
		{\True $24\,480$ cách} 
		 {$720$ cách}
		{$2\,520$  cách}
	\loigiai{ 
		Tìm bài toán đối đó là tìm số cách sao cho sau khi tặng xong có $1$  loại phần thưởng là hết.\\
		Trường  hợp 1: Hết cặp sách.\\
		Số cách chọn $4$  cặp sách là $1$ cách.\\
		Số cách chọn $1$ trong 6 phần thưởng  còn lại: $6$ cách.\\
		Số cách tặng $5$ phần thưởng đó cho $5$ em học  sinh là $\mathrm{A}_5^5=120$ cách.\\
		Vậy có $6\cdot 120=720$  cách.\\
		Trường  hợp 2: Hết hộp bút.\\
		Số cách  chọn $3$ hộp bút là $1$ cách.\\
		Số cách chọn  $2$ trong 7 phần thưởng còn lại là $\mathrm{C}_7^2$ cách.\\
		 Vậy có $21$ cách chọn sách.\\
		Số cách tặng $5$ phần  thưởng đó cho $5$ em học sinh là $\mathrm{A}_5^5=120$ cách.\\
		Vậy có $21\cdot  120=2\,520$ cách.\\
		Trường hợp 3: Hết vở. Tương tự  trường hợp $2$ thì có $2\,520$ cách.\\
		Số cách chọn $5$  phần thưởng trong $10$ phần thưởng và tặng cho $5$ em là $\mathrm{C}_{10}^5\cdot \mathrm{A}_5^5=30\,240$ cách.\\
		Vậy số cách chọn sao cho sau khi tặng xong, mỗi một trong ba loại phần thưởng trên đều còn ít nhất một phần thưởng là $30\,240-720-2\,520-2\,520=24\,480$  cách.
	 }
\end{ex}


%%%%%%%%%%%%%%%% ĐÚNG/SAI %%%%%%%%%%%%%%%%

% nguồn: MaHoaNH_DungSai_Lop10
\begin{ex}%[0D8H2-5]%[ma:bb14553b]
	Nhân dịp lễ sơ kết học kì I,cô An đã mua $10$ cuốn sách khác nhau để phát thưởng. Chọn khẳng định đúng trong các khẳng định  sau.
	\choiceTF 
		{\True Có $120$ cách chọn $3$ quyển sách để  phát thưởng}
		{\True có $720$ cách chọn $3$ quyển sách để phát thưởng cho $3$ học sinh, mỗi em nhận $1$ cuốn } 
		{Có $720$ cách chọn $3$ quyển sách để phát thưởng} 
		{Có $120$ cách chọn $3$ quyển sách để phát thưởng cho  $3$ học sinh, mỗi em nhận $1$ cuốn }
	\loigiai{ 
		\begin{itemize}
			\item Chọn $3$ quyển sách để phát thưởng có  $\mathrm{C}_{10}^3=120$ cách chọn.
			\item Chọn $3$ quyển sách để phát thưởng cho $3$  học sinh, mỗi em nhận $1$ cuốn có $\mathrm{A}_{10}^3=720$ cách chọn.
	\end{itemize}}
\end{ex}

% nguồn: MaHoaNH_DungSai_Lop10
\begin{ex}%[0D8H2-5]%[ma:73242a62]
	Một hộp đựng $2$ viên bi màu vàng và $3$ viên bi màu đỏ. Lấy $2$ viên bi từ hộp. Chọn khẳng định đúng trong các khẳng  định sau.
	\choiceTF 
		{\True  Có $10$ cách lấy ra $2$ viên bi trong hộp}
		{\True Có $6$ cách lấy ra $2$  viên bi khác màu}
		{\True Có $1$ cách lấy ra $2$ viên  bi màu vàng}
		{\True  Có $3$ cách lấy ra $2$ viên bi màu đỏ}
	 \loigiai{
		\begin{itemchoice}
			\itemch Lấy ra $2$ viên bi trong hộp có  $\mathrm{C}_5^2=10$ cách.
			\itemch Lấy ra $2$ viên bi khác  màu có $2\cdot 3=6$ cách.
			\itemch Lấy ra $2$ viên  bi màu vàng có $\mathrm{C}_2^2=1$ cách.
			\itemch Lấy ra $2$ viên bi màu đỏ có  $\mathrm{C}_3^2=3$ cách.
		\end{itemchoice}
	} 
\end{ex}

% nguồn: MaHoaNH_DungSai_Lop10
\begin{ex}%[0D8H2-5]%[ma:44089c63]
	Một hộp có $6$ viên bi xanh, $5$ viên bi đỏ và $4$ viên bi vàng,  chọn ngẫu nhiên $4$ viên bi. Khi đó, xét tính đúng sai của các mệnh đề sau.
	\choiceTF 
		{\True Chọn  $2$ bi xanh, $1$ bi đỏ và $1$ bi vàng: có $300$ cách}
		{ Chọn $1$ bi xanh, $2$ bi  đỏ và $1$ bi vàng: có $120$ cách}
		{\True Chọn $1$ bi xanh, $1$ bi đỏ và $2$ bi vàng:  có $180$ cách}
		{Có $600$ cách chọn ngẫu nhiên  $4$ viên bi từ hộp sao cho có đủ cả ba màu}
	\loigiai{ 
		\begin{itemchoice}
			\itemch Đúng. Chọn $2$ bi xanh, $1$ bi đỏ và $1$ bi vàng: có $\mathrm{C}_6^2 \cdot 5\cdot  4=300$ cách.
			\itemch Sai. Chọn $1$ bi xanh, $2$  bi đỏ và $1$ bi vàng: có $6\cdot \mathrm{C}_5^2\cdot 4=240$ cách.
			\itemch Đúng. Chọn $1$ bi xanh, $1$ bi đỏ và  $2$ bi vàng: có $6\cdot 5 \cdot \mathrm{C}_4^2=180$ cách.
			\itemch Sai. Số cách chọn ngẫu nhiên $4$ viên bi từ hộp sao cho có đủ cả  ba màu là $300+240+180=720$ cách.
		\end{itemchoice}
	} 
\end{ex}

% nguồn: MaHoaNH_DungSai_Lop10
\begin{ex}%[0D8H2-5]%[ma:58e5de49]
	Trong một hộp có $10$ quả cầu đỏ, $7$ quả cầu xanh. Lấy ra $3$ quả  cầu
	 \choiceTF
		{\True Có $\mathrm{C}_{10}^{3}$ cách lấy ra $3$ quả  cầu đỏ}
		{Có $\mathrm{A}_{7}^{3}$ cách  lấy ra $3$ quả cầu xanh}
		{\True Có $\mathrm{C}_{10}^{3}+\mathrm{C}_{7}^{3}$ cách lấy ra $3$ quả cầu  cùng màu}
		{Có $\mathrm{C}_{17}^{3}-\left(\mathrm{C}_{10}^{1}\mathrm{C}_{7}^{2}+\mathrm{C}_{10}^{2}\mathrm{C}_{7}^{1}\right)$ cách  lấy ra $3$ quả cầu khác màu}
	\loigiai{ 
		\begin{itemchoice}
			\itemch Lấy ra $3$ quả  cầu đỏ trong $10$ quả đỏ là $\mathrm{C}_{10}^{3}$.
			\itemch Lấy ra $3$ quả  cầu xanh trong $7$ quả đỏ là $\mathrm{C}_{7}^{3}$.
			\itemch Lấy ra $3$ quả cầu cùng màu có thể là $3$ quả  cầu đỏ hoặc $3$ quả cầu xanh nên có $\mathrm{C}_{10}^{3}+\mathrm{C}_{7}^{3}$.
			\itemch Lấy ra $3$ quả cầu khác màu có thể là $1$ quả cầu đỏ và $2$ quả cầu xanh hoặc $2$ quả  cầu đỏ và $1$ quả cầu xanh.\\
			Nên ta  có $\left(\mathrm{C}_{10}^{1}\mathrm{C}_{7}^{2}+\mathrm{C}_{10}^{2}\mathrm{C}_{7}^{1}\right)$.
		\end{itemchoice}
	 }
\end{ex}

% nguồn: MaHoaNH_DungSai_Lop10
\begin{ex}%[0D8H2-5]%[ma:b9e1d2b9]
	Một hộp có $6$ viên bi xanh, $5$ viên bi đỏ và $4$ viên bi vàng, chọn ngẫu nhiên $4$ viên bi. 
	 \choiceTF
		{\True Chọn $2$  bi xanh, $1$ bi đỏ và $1$ bi vàng có $300$ cách}
		{Chọn $1$ bi xanh, $2$ bi đỏ và $1$ bi vàng có  $120$ cách}
		{\True Chọn $1$ bi xanh, $1$ bi đỏ và $2$ bi vàng  có $180$ cách}
		{Có $600$ cách chọn ngẫu nhiên $4$ viên bi từ hộp sao cho có đủ cả ba  màu}
	\loigiai{ 
		Xét phép thử: chọn ngẫu nhiên $4$ viên bi từ một hộp có $6$ viên bi xanh, $5$ viên bi đỏ  và $4$ viên bi vàng.\\
		\begin{itemchoice}
			\itemch Số  cách chọn được $2$ bi xanh, $1$ bi đỏ và $1$ bi vàng là $\mathrm{C}_{6}^{2}\cdot 5\cdot 4=300$.
			\itemch Số cách chọn được $1$ bi xanh, $2$ bi đỏ và $1$ bi vàng là $6\cdot \mathrm{C}_{5}^{2}\cdot  4=240$.
			\itemch  Số cách chọn được $1$ bi xanh, $1$ bi đỏ và $2$ bi vàng là $6\cdot 5\cdot \mathrm{C}_{4}^{2}=180$.
			\itemch Để chọn được $4$ viên bi từ hộp sao cho có đủ cả  ba màu gồm $3$ trường hợp
			\begin{itemize}
				\item Chọn  $2$ bi xanh, $1$ bi đỏ và $1$ bi vàng. Có $300$ cách.
				\item Chọn $1$ bi xanh, $2$ bi đỏ  và $1$ bi vàng. Có $240$ cách.
				\item Chọn $1$ bi xanh, $1$ bi  đỏ và $2$ bi vàng. Có $180$ cách.
			\end{itemize}
			Vậy có  tất cả $720$ cách chọn $4$ viên bi có đủ cả ba màu.
		\end{itemchoice}
	 }
\end{ex}

% nguồn: MaHoaNH_DungSai_Lop10
\begin{ex}%[0D8H2-5]%[ma:c3760253]
	Trên một giá sách có $6$ quyển sách Toán, $5$ quyển sách Ngữ văn và $4$ quyển sách Tiếng Anh. Chọn ra  $4$ quyển sách.
	\choiceTF 
		{\True Số cách chọn ra  $4$ quyển sách bất kỳ là $1365$ cách}
		{\True Số cách chọn ra $4$ quyển sách  sao cho có đủ $3$ môn học là $720$ cách}
		{Số cách chọn ra $4$ quyển sách chỉ  có đúng $1$ môn học là $20$ cách}
		{Số cách  chọn ra $4$ quyển sách thuộc không quá $2$ môn học là $456$ cách}
	\loigiai{ 
		\begin{itemize}
			\item Mỗi cách chọn ra $4$ quyển  sách bất kỳ là một tổ hợp chập $4$ của $15$ phần tử.\\
			Vậy số cách chọn ra $4$ quyển sách bất  kỳ là $5\cdot\mathrm{C}_{15}^4=1365$ cách.
			\item Để $4$ quyển sách được chọn có đủ $3$ môn học,  có các trường hợp sau
			\begin{itemize}
				\item Trường hợp 1: $2$ sách Toán, $1$ sách Văn, $1$ sách Tiếng Anh,  có $\mathrm{C}_6^2 \cdot \mathrm{C}_5^1 \cdot \mathrm{C}_4^1= 300$ (cách).
				\item Trường hợp 2: $1$ sách Toán, $2$ sách Văn, $1$ sách Tiếng Anh, có $\mathrm{C}_6^1  \cdot \mathrm{C}_5^2 \cdot \mathrm{C}_4^1= 240$ (cách).
				\item Trường hợp 3: $1$ sách Toán, $1$ sách Văn, $2$ sách Tiếng  Anh, có $\mathrm{C}_6^1 \cdot \mathrm{C}_5^1 \cdot \mathrm{C}_4^2= 180$ (cách).
			\end{itemize}
			Vậy số cách chọn ra $4$ quyển sách sao cho có  đủ $3$ môn học là $300+240+180=720$ cách.
			\item Để chọn ra $4$ quyển sách chỉ có đúng $1$  môn học, ta có các trường hợp sau
			\begin{itemize}
				\item Chọn $4$ quyển  Toán có $\mathrm{C}_{6}^4=15$ cách.
				 \item Chọn $4$ quyển Ngữ Văn có $\mathrm{C}_{5}^4=5$ cách.
				\item Chọn $4$ quyển tiếng Anh  có $\mathrm{C}_{4}^4=1$ cách.
			\end{itemize}
			Vậy số cách chọn ra $4$  quyển sách chỉ có đúng $1$ môn học là $15+5+1=21$ cách.
			\item Để $4$ quyển sách được chọn thuộc không quá $2$ môn  học, xét hai trường hợp
			\begin{itemize}
				\item Trường hợp 1. Chọn $4$ quyển sách chỉ có 1 môn,  có $\mathrm{C}_4^4+\mathrm{C}_5^4+\mathrm{C}_6^4=21$ cách.
				\item Trường hợp 2.  Chọn $4$ quyển sách có $2$ môn. Có các khả năng sau
				\begin{itemize}
					\item Chọn Toán và Ngữ Văn, có $\mathrm{C}_{11}^4-\mathrm{C}_6^4-\mathrm{C}_5^4=310$  cách.
					\item Chọn  Ngữ Văn và tiếng Anh, có $\mathrm{C}_{9}^4-\mathrm{C}_4^4-\mathrm{C}_5^4=120$ cách.
					\item Chọn Toán và  tiếng Anh, có $\mathrm{C}_{10}^4-\mathrm{C}_6^4-\mathrm{C}_4^4=194$ cách.
				\end{itemize}
			\end{itemize}
			Vậy số cách chọn ra $4$ quyển sách thuộc không quá $2$ môn học là $21+310+120+194=645$  cách.
		\end{itemize}
		\begin{itemchoice}
			 \itemch Đúng.
			 \itemch Đúng.
			 \itemch Sai.
			 \itemch Sai.
		\end{itemchoice}
	} 
\end{ex}

% nguồn: MaHoaNH_DungSai_Lop10
\begin{ex}%[0D8V2-5]%[ma:a6b7d671]
	Từ một hộp chứa $ 12$ quả cầu trong đó có $ 8 $ quả màu đỏ, $ 3 $ quả màu xanh và $ 1  $ quả màu vàng, lấy ngẫu nhiên $ 3 $ quả. Xét tính đúng sai các khẳng định sau.
	\choiceTF{ Số cách  chọn ra $ 3 $ quả cầu từ hộp là $ 792$}
		{ Số cách chọn ra $ 3  $ quả cầu có đủ ba màu là $ 36$}
		{ Số  cách chọn ra $ 3 $ quả chỉ có một màu là $ 108$}
		{\True Số cách  chọn ra $ 3 $ quả có đúng hai màu là $ 139$}
	 \loigiai{
		\begin{itemchoice}
			\itemch Số cách chọn ra $ 3 $ quả cầu từ hộp  là $ \mathrm{C}_{12}^3=220 $ cách.
			\itemch Chọn ra $ 3 $ quả cầu có đủ ba màu bằng cách: Chọn $ 1 $ quả màu đỏ từ $ 8 $ quả màu đỏ, $ 1 $ quả màu xanh  từ $ 3 $ quả màu xanh và $ 1 $ quả màu vàng từ $ 1 $ quả màu vàng có $\mathrm{C}_8^1\cdot\mathrm{C}_3^1\cdot\mathrm{C}_1^1=24 $ cách.
			 \itemch Số cách chọn ra $ 3 $ quả chỉ có một màu là\\
			\textbf{TH1:} Chọn $ 3 $ quả đều là màu đỏ từ $ 8 $ quả màu đỏ có  $ \mathrm{C}_8^3 $ cách.\\
			\textbf{TH2:} Chọn $ 3 $ quả đều là màu xanh từ $  3 $ quả màu xanh có $ \mathrm{C}_3^3 $ cách.\\
			Theo quy tắc cộng số  cách chọn ra $ 3 $ quả chỉ có một màu là $ \mathrm{C}_8^3+\mathrm{C}_3^3=57$ cách.
			\itemch Số cách  chọn ra $ 3 $ quả có đúng hai màu là $ \mathrm{C}_{12}^3-\mathrm{C}_8^1\cdot\mathrm{C}_3^1\cdot\mathrm{C}_1^1-(\mathrm{C}_8^3+\mathrm{C}_3^3)=220-24-57=139 $ cách.
		\end{itemchoice}
	} 
\end{ex}

% nguồn: MaHoaNH_DungSai_Lop10
\begin{ex}%[0D8V2-5]%[ma:d6b96e3a]
	Trên một giá  sách có $6$ quyển sách Toán, $5$ quyển sách Ngữ văn và $4$ quyển sách Tiếng Anh. Chọn ra $4$ quyển sách. Xét tính Đúng - Sai của các khẳng định sau
	 \choiceTF
		{\True Số cách chọn ra $4$ quyển sách bất kỳ là  $1365$}
		{\True Số cách chọn ra $4$  quyển sách sao cho có đủ $3$ môn học là $720$}
		{Số cách chọn ra $4$ quyển  sách chỉ có đúng $1$ môn học là $20$}
		{Số cách  chọn ra $4$ quyển sách thuộc không quá $2$ môn học là $456$}
	 \loigiai{
		\begin{itemchoice}
			\itemch Số cách chọn ra $4$ quyển sách bất kỳ là số tổ hợp chập $4$ của $15$ phần  tử.\\
			Vậy số cách chọn ra $4$ quyển sách bất kỳ là  $5\cdot\mathrm{C}_{15}^4=1365$ cách.
			\itemch Ta có 
			\begin{itemize}
				\item Chọn $1$ quyển Toán, $1$ quyển Ngữ Văn và $2$  quyển tiếng Anh, có $\mathrm{C}_{6}^1\cdot\mathrm{C}_5^1\cdot\mathrm{C}_4^2=180$ cách.
				\item Chọn $1$ quyển Toán, $2$ quyển Ngữ Văn và $1$ quyển tiếng Anh,  có $\mathrm{C}_{6}^1\cdot\mathrm{C}_5^2\cdot\mathrm{C}_4^1=240$ cách.
				\item Chọn $2$ quyển Toán, $1$ quyển Ngữ Văn và $1$ quyển  tiếng Anh, có $\mathrm{C}_{6}^2\cdot\mathrm{C}_5^1\cdot\mathrm{C}_4^1=300$ cách.
			\end{itemize}
			Vậy số cách chọn ra $4$ quyển sách sao cho có  đủ $3$ môn học là $180+240+300=720$ cách.
			 \itemch Ta có
			\begin{itemize}
				\item  Chọn $4$ quyển Toán có $\mathrm{C}_{6}^4=15$ cách.
				 \item Chọn $4$ quyển Ngữ Văn có $\mathrm{C}_{5}^4=5$ cách.
				\item Chọn  $4$ quyển tiếng Anh có $\mathrm{C}_{4}^4=1$ cách.
			\end{itemize}
			Vậy số  cách chọn ra $4$ quyển sách chỉ có đúng $1$ môn học là $15+5+1=21$ cách.
			\itemch \textbf{Trường hợp 1.} Chọn $4$  quyển sách chỉ có 1 môn, có $\mathrm{C}_4^4+\mathrm{C}_5^4+\mathrm{C}_6^4=21$ cách.\\
			\textbf{Trường hợp 2.} Chọn $4$ quyển sách có $2$  môn.
			\begin{itemize}
				\item Chọn Toán và Ngữ  Văn, có $\mathrm{C}_{11}^4-\mathrm{C}_6^4-\mathrm{C}_5^4=310$ cách.
				\item Chọn  Ngữ Văn và tiếng Anh, có $\mathrm{C}_{9}^4-\mathrm{C}_4^4-\mathrm{C}_5^4=120$ cách.
				\item Chọn Toán và tiếng Anh,  có $\mathrm{C}_{10}^4-\mathrm{C}_6^4-\mathrm{C}_4^4=194$ cách.
			\end{itemize}
			Vậy số cách chọn ra $4$ quyển sách thuộc không quá $2$ môn học là  $21+310+120+194=645$ cách.
		\end{itemchoice}
	} 
\end{ex}


%%%%%%%%%%%%%%%% TRẢ LỜI NGẮN %%%%%%%%%%%%%%%%

% nguồn: MaHoaNH_TraLoiNgan_Lop10
\begin{bt}%[0D8H2-5]%[ma:0bc50c70]
	Phòng tranh của một họa sĩ có $12$ bức tranh. Họa sĩ đó có bao nhiêu cách chọn ra $3$ bức tranh trong đó để tham dự triển  lãm nghệ thuật?
	 \shortans{220}
	\loigiai{ 
		Vì các bức tranh chọn ra không sắp xếp  thứ tự nên mỗi cách chọn $3$ bức tranh trong $12$ bức tranh là một tổ hợp chập $3$ của $12$.\\
		Do đó  số cách chọn là $\mathrm{C}_{12}^{3}=220$.
	 }
\end{bt}

% nguồn: MaHoaNH_TraLoiNgan_Lop10
\begin{bt}%[0D8H2-5]%[ma:28e3c23b]
	Một hộp đựng $10$ viên bi có cùng kích thước, trong đó có $7$ viên bi màu đỏ và $3$ viên bi màu xanh. Hỏi có bao nhiêu cách chọn $2$  viên bi sao cho có ít nhất một viên bi màu xanh?
	\shortans{24} 
	\loigiai{ 
		Để chọn  $2$ viên bi có ít nhất một viên màu xanh, ta có các trường hợp
		\begin{itemize}
			\item Trường hợp 1: $2$ viên đều màu xanh,  có $\mathrm{C}_3^2=3$ cách.
			\item Trường  hợp 2: có một viên đỏ, một viên xanh, có $\mathrm{C}_3^1\cdot \mathrm{C}_7^1=21$ cách.
		\end{itemize}
		Theo quy tắc cộng,  có $3+21=24$ cách.
	 }
\end{bt}

% nguồn: MaHoaNH_TraLoiNgan_Lop10
\begin{bt}%[0D8H2-5]%[ma:21ba4f48]
	Bình $A$ chứa $3$ quả  cầu xanh, $4$ quả cầu đỏ và $5$ quả cầu trắng. Bình $B$ chứa $4$ quả cầu xanh, $3$ quả cầu đỏ và $6$ quả cầu trắng. Bình C chứa $5$ quả cầu xanh, $5$ quả cầu đỏ và $2$ quả cầu trắng. Từ mỗi bình lấy ra một quả cầu. Có bao nhiêu cách lấy để cuối cùng được $3$ quả có màu giống nhau.
	\shortans{180} 
	\loigiai{ 
		\begin{itemize}
			\item \textbf{Trường hợp 1:} Lấy được $3$ quả  cầu xanh từ $3$ bình\\
			Số cách  lấy: $\mathrm{C}_{3}^{1}\mathrm{C}_{4}^{1}\mathrm{C}_{5}^{1}=60$.
			\item \textbf{Trường hợp 2:} Lấy được $3$ quả cầu  đỏ từ $3$ bình\\
			Số cách  lấy: $\mathrm{C}_{4}^{1}\mathrm{C}_{3}^{1}\mathrm{C}_{5}^{1}=60$.
			 \item \textbf{Trường hợp 3:} Lấy được $3$ quả cầu trắng từ $3$ bình\\
			 Số cách lấy: $\mathrm{C}_{5}^{1}\mathrm{C}_{6}^{1}\mathrm{C}_{2}^{1}=60$.
		\end{itemize}
		Vậy có $60\cdot3=180$ cách lấy được $3$ quả cùng  màu từ $3$ bình.
	 }
\end{bt}

% nguồn: MaHoaNH_TraLoiNgan_Lop10
\begin{bt}%[0D8V2-5]%[ma:c179b341]
	Có $4$ quyển sách toán khác nhau,  $3$ quyển sách lý khác nhau, $2$ quyển sách hoá khác nhau. Có bao nhiêu cách xếp số sách đó trên một giá nằm ngang sao cho các sách cùng loại nằm cạnh nhau.
	\shortans{1728} 
	 \loigiai{
		Đề xếp những quyển sách trên theo yêu cầu bài toán  ta tiến hành các giai đoạn sau\\
		Giai đoạn $1$: Xếp $4$ quyển  sách toán ta có $4!$ cách.\\
		Giai đoạn $2$: Xếp  $3$ quyển sách lý ta có $3!$ cách.\\
		Giai đoạn $3$: Xếp $2$  quyển sách hoá ta có $2!$ cách.\\
		Giai đoạn $4$: Xem các sách cùng loại toán, lý,  hoá lần lượt là các tập hợp: T, L, H. Xếp ba tập hợp này ta có $3!$ cách.\\
		Vậy theo quy tắc nhân ta có $4!\cdot 3!\cdot 2!\cdot  3!=1728$ cách.
	 }
\end{bt}

% nguồn: MaHoaNH_TraLoiNgan_Lop10
\begin{bt}%[0D8V2-5]%[ma:8bdf067f]
	Một người chọn ngẫu nhiên $4$ chiếc giày từ $5$ đôi giày cỡ khác nhau. Tính số cách chọn sao cho trong $4$ chiếc giày được  chọn không có $2$ chiếc nào tạo thành một đôi.
	\shortans{80} 
	 \loigiai{
		 Ta có
		\begin{itemize}
			\item Chọn  $4$ đôi trong $5$ đôi có $\mathrm{C}_{5}^{4}$ cách.
			\item Trong mỗi đôi được chọn lấy một chiếc: có $2\cdot 2\cdot  2\cdot 2=16$ cách.
		\end{itemize}
		Do đó có  $\mathrm{C}_{5}^{4}\cdot 16=80$ cách chọn.
	} 
\end{bt}

% nguồn: MaHoaNH_TraLoiNgan_Lop10
\begin{bt}%[0D8V2-5]%[ma:0f8a0a04]
	Một hộp chứa $5$ quả cầu trắng và $6$ quả cầu đỏ, các quả cầu có kích thước và khối lượng giống nhau. Có bao nhiêu cách lấy ngẫu nhiên đồng thời $3$ quả cầu trong hộp sao cho trong $3$ quả cầu lấy ra có ít nhất một  quả cầu trắng?
	 \shortans{145}
	\loigiai{ 
		Xét  các trường hợp sau
		\begin{itemize}
			 \item Trường hợp 1: lấy được $1$ quả cầu trắng và $2$ quả cầu đỏ.\\
			Số cách lấy $1$  quả cầu trắng và $2$ quả cầu đỏ là $\mathrm{C}_5^1 \cdot \mathrm{C}_6^2=75$.
			\item Trường hợp 2: lấy được $2$ quả cầu trắng và $1$ quả  cầu đỏ.\\
			Số cách lấy $1$ quả cầu trắng  và $2$ quả cầu đỏ là $\mathrm{C}_5^2 \cdot \mathrm{C}_6^1=60$.
			\item Trường hợp 3: lấy được $3$ quả cầu  trắng. \\
			 Số cách lấy $3$ quả cầu trắng là $\mathrm{C}_5^3= 10$.
		\end{itemize}
		Vậy số cách lấy ngẫu nhiên đồng thời $3$ quả cầu trong hộp sao cho trong $3$ quả cầu lấy ra  có ít nhất một quả cầu trắng là
		 \[
		75+60+10=145. 
		\] 
	} 
\end{bt}


%%%%%%%%%%%%%%%% TỰ LUẬN %%%%%%%%%%%%%%%%

% nguồn: MaHoaNH_TuLuan_Lop10
\begin{bt}%[0D8H2-5]%[ma:ef273c3a]
	Có bao nhiêu cách xếp $6$ quyển sách Toán và $6$ quyển sách Tiếng Anh (các quyển sách là khác nhau) vào một hàng ngang của giá sách  sao cho các quyển sách Toán và sách Tiếng Anh ở vị trí xen kẽ nhau?
	\loigiai{ 
		Giả  sử trên hàng ngang của giá sách có đánh số từ $1$ đến $12$.\\
		Nếu $6$ quyển sách Toán được xếp vào vị trí lẻ thì $6$ quyển sách Tiếng Anh xếp vào các vị trí còn lại nên có $6!$ cách xếp quyển sách Toán và $6!$ cách xếp quyển sách Tiếng Anh. Suy ra có $(6!)^2=518\,400$ cách  xếp.\\
		Nếu $6$ quyển sách Toán được xếp vào vị trí chẵn thì $6$ quyển sách Tiếng Anh xếp vào các vị trí còn lại nên có $6!$ cách xếp quyển sách Toán  và $6!$ cách xếp quyển sách Tiếng Anh. Suy ra có $(6!)^2=518\,400$ cách xếp.\\
		Vậy có tất cả $2\cdot  518\,400=1\,036\,800$ cách sắp xếp sao cho các quyển sách Toán và sách Tiếng Anh ở vị trí xen kẽ nhau.
	 }
\end{bt}

% nguồn: MaHoaNH_TuLuan_Lop10
\begin{bt}%[0D8H2-5]%[ma:8d8106d1]
	Từ một bó gồm $5$ bông hoa đỏ, $6$ bông hoa vàng, $7$ bông hoa tím. Có bao nhiêu cách chọn ra $4$ bông hoa có đủ  cả $3$ màu.
	\loigiai{ 
		Trường hợp 1: $2$  hoa đỏ, $1$ hoa vàng, $1$ hoa tím:
		 \[n_1 = \mathrm{C}_5^2 \cdot \mathrm{C}_6^1 \cdot \mathrm{C}_7^1 = 420 \text{ cách.} \]
		Trường hợp 2: $1$  hoa đỏ, $2$ hoa vàng, $1$ hoa tím:
		\[n_2 = \mathrm{C}_5^1 \cdot \mathrm{C}_6^2 \cdot \mathrm{C}_7^1 = 525 \text{  cách.} \]
		Trường hợp 3: $1$ hoa  đỏ, $1$ hoa vàng, $2$ hoa tím:
		\[n_3 = \mathrm{C}_5^1  \cdot \mathrm{C}_6^1 \cdot \mathrm{C}_7^2 = 630 \text{ cách.} \]
		Vậy số cách chọn $4$ bông hoa có  đủ $3$ màu là:
		 \[n = n_1 + n_2 + n_3 = 420 + 525 + 630 = 1575 \text{ cách.} \]
	 }
\end{bt}

% nguồn: MaHoaNH_TuLuan_Lop10
\begin{bt}%[0D8H2-5]%[ma:ab48d1c9]
	Trong một môn học, thầy giáo có $30$ câu hỏi khác nhau gồm $5$ câu hỏi khó, $10$ câu hỏi trung bình, $15$ câu hỏi dễ. Từ $30$ câu hỏi đó có thể lập được bao nhiêu đề để kiểm tra, mỗi đề gồm $5$ câu hỏi khác nhau, sao cho mỗi đề thi nhất thiết phải có đủ $3$ loại (khó, trung bình, dễ)  và số câu hỏi dễ không ít hơn $2$?
	 \loigiai{
		Mỗi đề kiểm tra phải có số câu dễ là $2$ hoặc $3$, nên ta có  các trường hợp sau:\\
		Đề có $2$ câu dễ, $2$ câu trung bình, $1$ câu khó, thì có  số cách chọn là:
		\[\mathrm{C}_{15}^2 \cdot \mathrm{C}_{10}^2  \cdot \mathrm{C}_5^1 = 23625.\]
		 Đề có $2$ câu dễ, $1$ câu trung bình, $2$ câu khó, thì có số cách chọn là:
		 \[\mathrm{C}_{15}^2 \cdot \mathrm{C}_{10}^1 \cdot \mathrm{C}_5^2 = 10500. \]
		Đề có $3$ câu dễ, $1$ câu  trung bình, $1$ câu khó, thì có số cách chọn là:
		\[\mathrm{C}_{15}^3 \cdot \mathrm{C}_{10}^1 \cdot  \mathrm{C}_5^1 = 22750.\]
		Vì các cách chọn trên đôi một khác nhau, nên số đề kiểm tra có thể lập được là: 
		 \[23625 + 10500 + 22750 = 56875.\]
	} 
\end{bt}

% nguồn: MaHoaNH_TuLuan_Lop10
\begin{bt}%[0D8H2-5]%[ma:d6f44e60]
	Có $5$  bông hồng, $4$ bông trắng (mỗi bông đều khác nhau về hình dáng). Một người cần chọn một bó gồm $4$ bông từ số bông này.
	\begin{enumerate}
		 \item Có bao nhiêu cách chọn sao cho số bông mỗi màu bằng nhau?
		\item Có bao nhiêu cách chọn sao cho có đủ cả hai  màu?
	\end{enumerate}
	 \loigiai{
		\begin{listEX}[1]
			\item Số cách chọn $2$  bông hồng từ $5$ bông hồng là $\mathrm{C}_5^2$ (cách).\\
			Số cách chọn $2$ bông trắng từ  $4$ bông trắng là $\mathrm{C}_4^2$ (cách).\\
			Số cách chọn một bó bông thỏa  mãn đề bài là $\mathrm{C}_5^2 \cdot \mathrm{C}_4^2=60$ (cách).
			 \item
			{\bfseries Trường hợp 1.} Chọn $3$ bông hồng, $1$ bông  trắng: có $\mathrm{C}_5^3 \cdot \mathrm{C}_4^1=40$ (cách).\\
			{\bfseries Trường hợp 2.} Chọn $2$ bông hồng, $2$  bông trắng: có $\mathrm{C}_5^2 \cdot \mathrm{C}_4^2=60$ (cách).\\
			{\bfseries Trường hợp 3.} Chọn $1$  bông hồng, $3$ bông trắng: có $\mathrm{C}_5^1 \cdot \mathrm{C}_4^3=20$ (cách).\\
			Theo quy tắc cộng ta  có tất cả $40+60+20=120$ (cách chọn).
		\end{listEX}
	} 
\end{bt}

% nguồn: MaHoaNH_TuLuan_Lop10
\begin{bt}%[0D8H2-5]%[ma:dc1737de]
	Từ $5$  bông hồng vàng, $3$ bông hồng trắng, $4$ bông hồng đỏ (các bông hồng xem như đôi một khác nhau). Người ta muốn chọn ra $1$ bó hoa hồng gồm $7$ bông. Có bao nhiêu cách chọn?
	\begin{enumerate}
		\item 1 bó hoa trong đó có đúng một bông hồng  đỏ.
		\item 1 bó hoa trong đó có  ít nhất $3$ bông hồng vàng và ít nhất $3$ bông hồng đỏ.
	\end{enumerate}
	 \loigiai{
		\begin{enumerate}
			\item Chọn $1$ bó  hoa gồm $7$ bông, trong đó có đúng $1$ bông hồng đỏ, $6$ bông hồng còn lại chọn trong $8$ bông (gồm vàng và trắng). \\
			Số cách chọn $\mathrm{C}_4^1\cdot \mathrm{C}_8^6=112$  cách.
			\item Có các trường hợp sau xảy ra thỏa yêu cầu bài toán 
			\begin{itemize}
				\item Trường hợp 1: Chọn $3$ bông  hồng vàng, $3$ bông hồng đỏ và $1$ bông hồng trắng
				Khi  đó có $\mathrm{C}_5^3\cdot \mathrm{C}_4^3\cdot \mathrm{C}_3^1$ cách.
				\item Trường hợp 2: Chọn $4$ bông hồng  vàng và $3$ bông hồng đỏ có $\mathrm{C}_5^4\cdot \mathrm{C}_4^3$ cách.
				\item Trường hợp 3: Chọn $3$ bông hồng vàng và  $4$ bông hồng đỏ có $\mathrm{C}_5^3\cdot \mathrm{C}_4^4$ cách.
			\end{itemize}
			Theo quy tắc cộng có $\mathrm{C}_5^3\cdot \mathrm{C}_4^3\cdot \mathrm{C}_3^1+\mathrm{C}_5^4\cdot \mathrm{C}_4^3+\mathrm{C}_5^3\cdot \mathrm{C}_4^4$  cách.
		\end{enumerate}
	 }
\end{bt}

% nguồn: MaHoaNH_TuLuan_Lop10
\begin{bt}%[0D8H2-5]%[ma:8499877d]
	Có $9$ viên bi  xanh, $5$ viên bi đỏ, $4$ bi vàng có kích thước đôi một khác nhau.
	\begin{enumerate}
		\item Có bao nhiêu cách chọn ra $6$ viên bi, trong đó có  đúng $2$ viên bi đỏ.
		\item Có bao nhiêu cách chọn ra $6$ viên bi,  trong đó số bi xanh bằng số bi đỏ.
	\end{enumerate}
	\loigiai{ 
	\begin{enumerate}
		\item Ta  lần lượt thực hiện các công đoạn sau:\\
		Chọn $2$  bi đỏ trong $5$ bi đỏ, có $\mathrm{C}_5^2$ cách chọn. \\
		 Có $\mathrm{C}_{13}^4$ cách chọn $4$ bi trong $13$ viên bi xanh và vàng. \\
		Vậy ta có $\mathrm{C}_5^2\cdot \mathrm{C}_{13}^4=7\,150$ cách. 
		\item Số bi xanh, đỏ,  vàng được chọn có $3$ trường hợp là
		\begin{itemize}
			\item Trường hợp 1: Chọn $3$ xanh, $3$ đỏ,  ta có $\mathrm{C}_9^3\mathrm{C}_5^3$ cách.
			\item  Trường hợp 2: Chọn $2$ xanh, $2$ đỏ, $2$ vàng, ta có $\mathrm{C}_9^2\mathrm{C}_5^2\mathrm{C}_4^2$ cách.
			\item Trường hợp 3: Chọn $1$ xanh, $1$ đỏ, $4$ vàng, ta có $\mathrm{C}_9^1\mathrm{C}_5^1\mathrm{C}_4^4$  cách.
		\end{itemize}
		Theo quy tắc cộng ta có $\mathrm{C}_9^3\cdot \mathrm{C}_5^3+\mathrm{C}_9^2\cdot \mathrm{C}_5^2\cdot \mathrm{C}_4^2+\mathrm{C}_9^1\cdot \mathrm{C}_5^1\cdot  \mathrm{C}_4^4=3\,045$ cách.
	\end{enumerate}
	 }
\end{bt}

% nguồn: MaHoaNH_TuLuan_Lop10
\begin{bt}%[0D8H2-5]%[ma:e216c3db]
	Có một hộp đựng $5$ viên bi xanh, $6$  viên bi đỏ và $4$ viên bi vàng.
	\begin{enumerate}
		 \item Có bao nhiêu cách lấy ra $6$ viên bi, trong đó có $2$ viên bi xanh và có nhiều nhất $2$ viên bi vàng và phải có đủ $3$ màu.
		\item Có bao nhiêu cách lấy ra $9$  viên bi có đủ $3$ màu.
	\end{enumerate}
	\loigiai{ 
		\begin{enumerate}
			\item Các trường  hợp xảy ra theo yêu cầu đề:
			\begin{itemize}
				\item Trường hơp 1: $2$ xanh, $2$ vàng, $2$ đỏ, có $\mathrm{C}_5^2\cdot  \mathrm{C}_4^2\cdot \mathrm{C}_6^2$ cách.
				\item Trường hợp 2: $2$ xanh, $1$ vàng, $3$ đỏ, có $\mathrm{C}_5^2\cdot  \mathrm{C}_4^1\cdot \mathrm{C}_6^3$ cách.
			\end{itemize}
			Vậy có $\mathrm{C}_5^2\cdot \mathrm{C}_4^2\cdot \mathrm{C}_6^2+\mathrm{C}_5^2\cdot \mathrm{C}_4^1\cdot \mathrm{C}_6^3=1\,700$ cách. 
			\item Sử dụng phương  pháp gián tiếp: Lấy ra $9$ viên bi trong $15$ viên bi bất kỳ, có $\mathrm{C}_{15}^9$ cách.
			\begin{itemize}
				\item Trường hợp 1: lấy $9$ viên bi chỉ có $2$ màu là  xanh và đỏ, có $\mathrm{C}_{11}^9$ cách.
				\item Trường hợp 2: lấy $9$ viên bi chỉ có $2$ màu là xanh và vàng, có $\mathrm{C}_9^9$  cách.
				\item Trường hợp 3:  lấy $9$ viên bi chỉ có $2$ màu là đỏ và vàng, có $\mathrm{C}_{10}^9$ cách.
			\end{itemize}
			 Vậy có $\mathrm{C}_{15}^9-(\mathrm{C}_{11}^9+\mathrm{C}_9^9+\mathrm{C}_{10}^9)=4\,984$ cách.
		\end{enumerate}
	 }
\end{bt}
